\documentclass[12pt,a4paper]{article}
\usepackage[margin=2cm]{geometry}
\usepackage{amsmath,amssymb,amsthm,mathtools}
\usepackage{enumitem}
\usepackage{algorithm}
\usepackage{algpseudocode}
\usepackage{microtype}
\usepackage[hidelinks]{hyperref}

\setlist{nosep}
\setlength{\parindent}{0pt}
\setlength{\parskip}{0.4em}

\newtheorem{theorem}{Theorem}[section]
\newtheorem{lemma}[theorem]{Lemma}
\newtheorem{proposition}[theorem]{Proposition}
\newtheorem{corollary}[theorem]{Corollary}

\theoremstyle{definition}
\newtheorem{definition}[theorem]{Definition}
\newtheorem{example}[theorem]{Example}
\newtheorem{problem}{Problem}[section]

\theoremstyle{remark}
\newtheorem*{remark}{Remark}

\newenvironment{solution}{\begin{proof}[Solution]}{\end{proof}}

\DeclareMathOperator{\OPT}{OPT}
\DeclareMathOperator{\poly}{poly}

\DeclareMathOperator{\ppi}{ppi}


\title{Advanced Algorithms\\\large Klevis Imeri Lecture Notes}
\author{}
\date{Autumn 2026}

\begin{document}

\maketitle
\tableofcontents

\section{Approximation Algorithms}

\begin{definition}[Problem]
\[
\Pi:=(\mathcal I,\mathcal S).
\]
\[
\begin{aligned}
\text{Input:}\quad&I\in\mathcal I,\\
\text{Output:}\quad&S\in\mathcal S_I.
\end{aligned}
\]
\[
\begin{aligned}
\Pi&:\text{ problem},\\
\mathcal I&:\text{ instance space},\qquad I:\text{ instance},\\
\mathcal S&:\text{ set of possible outputs},\qquad S:\text{ output},\\
\mathcal S_I&:=\{\,S\in\mathcal S:S\text{ is valid for }I\,\}.
\end{aligned}
\]
\end{definition}

\noindent\begin{minipage}{\linewidth}
\begin{definition}[Optimization problem]
\[
\Pi:=(\mathcal I,\mathcal S).
\]
\[
c:\mathcal I\times\mathcal S\to\mathbb R_{\geq0}.
\]
\[
\text{For fixed }I:\quad c(S):=c(I,S),\qquad c(\OPT(I)):=c(I,\OPT(I)).
\]
\[
\OPT(I)\in\mathcal S_I:=
\begin{cases}
\displaystyle\operatorname*{arg\,min}_{T\in\mathcal S}c(T), & \text{minimization},\\[0.6em]
\displaystyle\operatorname*{arg\,max}_{T\in\mathcal S}c(T), & \text{maximization}.
\end{cases}
\quad\text{+ other constraints}.
\]
\[
\begin{aligned}
\Pi&:\text{ optimization problem},\\
\mathcal I&:\text{ instance space},\qquad I:\text{ instance},\\
\mathcal S&:\text{ set of possible outputs},\qquad S:\text{ output},\\
\mathcal S_I&:\text{ optimal solutions for }I,\\
c&:\text{ cost function},\qquad c(S):\text{ cost of }S\text{ for }I,\\
\OPT(I)&:\text{ optimal solution for }I\text{ (assumed to exist)}.
\end{aligned}
\]
\end{definition}
\end{minipage}\par

\begin{definition}[Algorithm]
\[
\mathcal A:\mathcal I\to\mathcal S,
\qquad
\forall I\in\mathcal I:\ \mathcal A(I)\in\mathcal S_I.
\]
\[
\Pi:=(\mathcal I,\mathcal S).
\]
\[
\begin{aligned}
\Pi&:\text{ problem},\\
\mathcal I&:\text{ instance space},\qquad I:\text{ instance},\\
\mathcal S&:\text{ set of possible outputs},\\
\mathcal S_I&:=\{\,S\in\mathcal S:S\text{ is valid for }I\,\},\\
\mathcal A&:\text{ algorithm},\qquad
\mathcal A(I):\text{ solution returned on }I.
\end{aligned}
\]
\end{definition}

\begin{definition}[Polynomial]
\[
\poly(n):=\{\,f:\exists c\in\mathbb N,\ f\in O(n^{c})\,\}.
\]
\[
\begin{aligned}
\poly(n)&:\text{ functions bounded by a polynomial in }n,\\
n&:\text{ input size},\\
f&:\text{ function},\\
c&:\text{ constant exponent},\\
O(n^c)&:\text{ asymptotic upper bound of order }n^c.
\end{aligned}
\]
\end{definition}

\begin{definition}[Running time (TODO)]
\end{definition}

\noindent\begin{minipage}{\linewidth}
\begin{definition}[Efficient algorithm]
\[
\exists p\in\poly(n)\ \forall I\in\mathcal I:
\quad \mathrm{time}_{\mathcal A}(I)\leq p(|I|).
\]
\[
\begin{aligned}
\mathcal A&:\text{ algorithm},\\
\mathcal I&:\text{ instance space},\qquad I:\text{ instance},\\
|I|&:\text{ encoding length of the instance }I,\\
\mathrm{time}_{\mathcal A}(I)&:\text{ running time of }\mathcal A\text{ on }I,\\
p&:\text{ polynomial bound},\qquad n:\text{ input size}.
\end{aligned}
\]
\end{definition}
\end{minipage}\par

\noindent\begin{minipage}{\linewidth}
\begin{definition}[$\alpha$-approximation]
\label{def:alpha-approximation}
\[
\begin{gathered}
\mathcal A\ \text{is an}\ \alpha\text{-approximation}
\iff\\[3pt]
\bigl[\exists p\in\poly(n)\ \forall I\in\mathcal I:
\mathrm{time}_{\mathcal A}(I)\leq p(|I|)\bigr]\ \land\\[3pt]
\forall I\in\mathcal I:\\[3pt]
\begin{cases}
c\bigl(\mathcal A(I)\bigr)\leq\alpha(I)\,c\bigl(\OPT(I)\bigr), & \text{minimization},\\[0.6em]
c\bigl(\OPT(I)\bigr)\leq\alpha(I)\,c\bigl(\mathcal A(I)\bigr), & \text{maximization}.
\end{cases}
\end{gathered}
\]
\[
\Pi:=(\mathcal I,\mathcal S),
\qquad \alpha:\mathcal I\to\mathbb R_{\geq1}.
\]
\[
\mathcal A:\mathcal I\to\mathcal S,
\qquad \forall I\in\mathcal I:\ \mathcal A(I)\text{ is feasible for }I.
\]
\[
\begin{aligned}
\Pi&:\text{ optimization problem},\\
\mathcal I&:\text{ instance space},\qquad I:\text{ instance},\\
\mathcal S&:\text{ set of possible outputs},\\
\mathcal S_I&:\text{ optimal solutions for }I,\\
c&:\text{ cost function},\\
\mathcal A&:\text{ algorithm},\\
\OPT(I)&:\text{ optimal solution for }I,\\
\alpha&:\text{ approximation-factor function},\\
\alpha(I)&:\text{ approximation factor for }I.
\end{aligned}
\]
\[
\text{Constant factor }a\geq1:\quad \alpha(I):=a\quad(\forall I\in\mathcal I).
\]
\end{definition}
\end{minipage}\par

\noindent\begin{minipage}{\linewidth}
\begin{definition}[Minimum set cover]
\[
\Pi_{\text{minimum set cover}}:=(\mathcal I,\mathcal S).
\]
\[
\begin{aligned}
\mathcal I&:=\left\{\,(U,S,w):
1\leq|U|<\infty,\ S\subseteq 2^U,\
\bigcup_{S_i\in S}S_i=U,\ w:S\to\mathbb R_{\geq 0}\,\right\},\\[3pt]
I&=(U,S,w)\in\mathcal I,\qquad c(T):=\sum_{S_i\in T}w(S_i),\\[3pt]
\OPT(I)\in\mathcal S_I&:=\left\{\,T\subseteq S:
\displaystyle\bigcup_{S_i\in T}S_i=U
\land\ T\in\operatorname*{arg\,min}_{\substack{T'\subseteq S\\\bigcup_{S_i\in T'}S_i=U}}c(T')
\,\right\}.
\end{aligned}
\]
\[
\begin{aligned}
\mathcal I&:\text{ instance space},\\
\mathcal S&:\text{ set of possible outputs},\\
\mathcal S_I&:\text{ minimum-cost covers of }U,\\
c&:\text{ cost function},\\
I&:\text{ instance},\\
U&:\text{ nonempty finite universe},\\
S&:\text{ collection of available subsets of }U,\\
S_i&:\text{ individual available set},\\
T&:\text{ selected subcollection},\\
w(S_i)&:\text{ cost of set }S_i,\\
c(T)&:\text{ total cost of the selected sets},\\
2^U&:\text{ power set of }U.
\end{aligned}
\]
\end{definition}
\end{minipage}\par

\noindent\begin{minipage}{\linewidth}
\begin{example}[Minimum set cover]
\[
U:=\{e_1,e_2,e_3,e_4,e_5\},
\qquad S:=\{S_1,S_2,S_3,S_4\},
\qquad I:=(U,S,w).
\]
\[
\begin{array}{c|c|c}
i & S_i & w(S_i):=i^2\\ \hline
1 & \{e_1,e_2,e_5\} & 1\\
2 & \{e_4\} & 4\\
3 & \{e_2,e_3\} & 9\\
4 & \{e_1,e_4,e_5\} & 16
\end{array}
\]
\[
T:=\{S_3,S_4\}\subseteq S,\qquad\bigcup_{S_i\in T}S_i=U,
\qquad c(T)=9+16=25.
\]
\[
T^*:=\OPT(I)=\{S_1,S_2,S_3\},
\qquad c(T^*)=1+4+9=14.
\]
\[
|T^*|=3>2=|T|,
\qquad c(T^*)=14<25=c(T).
\]
\end{example}
\end{minipage}\par

\noindent\begin{minipage}{\linewidth}
\begin{theorem}[Complexity of minimum set cover]
\[
\Pi_{\text{minimum set cover}}\in\mathrm{NP\text{-}hard}.
\]
\[
\mathrm{SC}_{\mathrm{dec}}:=\left\{\,(U,S,k):
\begin{array}{l}
|U|<\infty,\ S\subseteq 2^U,\ k\in\mathbb N,\\
\exists T\subseteq S:\ |T|\leq k,\ \displaystyle\bigcup_{S_i\in T}S_i=U
\end{array}\,\right\}.
\]
\[
\mathrm{SC}_{\mathrm{dec}}\in\mathrm{NPC}.
\]
\end{theorem}

\begin{proof}[Proof (Intuition)]
A proposed cover can be checked in polynomial time: count its sets and check
that their union is the universe. Thus the decision problem is in NP.

Vertex Cover, a known NP-complete problem, asks whether at most $k$ vertices
can touch every edge of a graph. Turn each edge into an element of $U$, and
each vertex into the set of edges touching it. Selecting vertices that touch
every edge is then equivalent to selecting sets that cover $U$. This
construction takes polynomial time, so Set Cover is NP-hard and hence
NP-complete as a decision problem.

With every set cost equal to $1$, an optimal cover also tells us whether a
cover with at most $k$ sets exists. Finding an optimal cover is therefore
NP-hard.
\end{proof}
\end{minipage}\par

\noindent\begin{minipage}{\linewidth}
\begin{definition}[Greedy set cover algorithm]
\[
\mathcal A(I):=\mathsf{GreedySetCover}(U,S,w).
\]
\[
\ppi(S_j):=\begin{cases}
\dfrac{w(S_j)}{|S_j|},&S_j\neq\varnothing,\\[3pt]
\infty,&S_j=\varnothing.
\end{cases}
\]
\begin{algorithm}[H]
\caption{$\mathsf{GreedySetCover}(U,S,w)$}
\begin{algorithmic}[1]
\If{$U=\varnothing$}
    \State \Return $\varnothing$
\Else
    \State $S_i\in\displaystyle\operatorname*{arg\,min}_{S_j\in S}\ppi(S_j)$
    \State $U'\gets U\setminus S_i$
    \State $S'\gets\{\,S_j':=S_j\setminus S_i:S_j\in S\,\}$
    \State $w(S_j')\gets w(S_j)\quad(\forall S_j'\in S')$
    \State \Return $\{S_i\}\cup\mathsf{GreedySetCover}(U',S',w)$
\EndIf
\end{algorithmic}
\end{algorithm}
\[
\Pi_{\text{minimum set cover}}=(\mathcal I,\mathcal S).
\]
\[
\begin{aligned}
U&:\text{ remaining universe},\\
S&:\text{ collection of remaining sets},\\
w(S_j)&:\text{ cost of set }S_j,\\
\ppi(S_j)&:\text{ cost per remaining element of }S_j,\\
S_i&:\text{ set chosen in the current call},\\
U'&:\text{ universe after removing the elements of }S_i,\\
S_j'&:\text{ remaining elements of }S_j,\\
S'&:\text{ collection of sets for the recursive call}.
\end{aligned}
\]
\end{definition}
\end{minipage}\par

\noindent\begin{minipage}{\linewidth}
\begin{example}[Greedy set cover]
\[
U_0:=\{e_1,e_2,e_3,e_4,e_5\},\qquad
S:=\{S_1,S_2,S_3,S_4\}.
\]
\[
\begin{aligned}
S_1&=\{e_1,e_2,e_5\},&\qquad S_2&=\{e_4\},\\
S_3&=\{e_2,e_3\},& S_4&=\{e_1,e_4,e_5\}.
\end{aligned}
\]
\[
w(S_j):=j^2\quad(j\in\{1,2,3,4\}),\qquad
I:=(U_0,S,w).
\]
\[
S_j^{(t)}:=S_j\cap U_t,\qquad
S^{(t)}:=\{\,S_j^{(t)}:j\in\{1,2,3,4\}\,\}.
\]
\[
\ppi_{j,t}:=\ppi(S_j^{(t)})=\begin{cases}
\dfrac{w(S_j)}{|S_j^{(t)}|},&S_j^{(t)}\neq\varnothing,\\[3pt]
\infty,&S_j^{(t)}=\varnothing.
\end{cases}
\]
\[
i_t\in\operatorname*{arg\,min}_{j\in\{1,2,3,4\}}\ppi_{j,t},\qquad
U_{t+1}:=U_t\setminus S_{i_t}.
\]
\[
\renewcommand{\arraystretch}{1.5}
\begin{array}{c|c|cccc|c|c}
t&U_t&\ppi_{1,t}&\ppi_{2,t}&\ppi_{3,t}&\ppi_{4,t}&S_{i_t}&U_{t+1}\\ \hline
0&U_0&\frac13&4&\frac92&\frac{16}{3}&S_1&\{e_3,e_4\}\\
1&\{e_3,e_4\}&\infty&4&9&16&S_2&\{e_3\}\\
2&\{e_3\}&\infty&\infty&9&\infty&S_3&\varnothing
\end{array}
\]
\[
\mathsf{GreedySetCover}(U_3,S^{(3)},w)=\varnothing.
\]
\[
\begin{aligned}
\mathcal A(I)
&=\mathsf{GreedySetCover}(U_0,S^{(0)},w)\\
&=\{S_1\}\cup\bigl(\{S_2\}\cup(\{S_3\}\cup\varnothing)\bigr)\\
&=\{S_1,S_2,S_3\},\\[3pt]
c\bigl(\mathcal A(I)\bigr)&=1+4+9=14.
\end{aligned}
\]
\end{example}
\end{minipage}\par

\noindent\begin{minipage}{\linewidth}
\begin{theorem}[Termination, feasibility, and efficiency of greedy set cover]
\label{thm:greedy-efficient}
\begin{flalign*}
&\forall I=(U,S,w)\in\mathcal I:\quad
\mathsf{GreedySetCover}(U,S,w)\text{ terminates, is feasible, and is efficient}.&&
\end{flalign*}
\end{theorem}
\begin{proof}
\setlength{\jot}{2pt}
\setlength{\abovedisplayskip}{3pt}
\setlength{\belowdisplayskip}{3pt}
\setlength{\abovedisplayshortskip}{3pt}
\setlength{\belowdisplayshortskip}{3pt}
\begin{flalign*}
&\begin{aligned}
&I=(U,S,w)\in\mathcal I,\qquad \mathcal A(I):=\mathsf{GreedySetCover}(U,S,w),\\
&U_t:\text{ remaining universe after }t\text{ selections},\qquad U_0:=U.
\end{aligned}&&
\end{flalign*}
\begin{flalign*}
&\text{1. Termination:}\quad
\mathcal A(I)\text{ terminates}\ \iff\ \exists r\in\mathbb N_{\geq0}:\ U_r=\varnothing.&&
\end{flalign*}
\begin{flalign*}
&\begin{aligned}
&i_t:\text{ index of the greedy set selected at step }t,\\
&U_t\neq\varnothing\ \Longrightarrow\ |S_{i_t}\cap U_t|\geq1,\\
&U_{t+1}:=U_t\setminus S_{i_t},\qquad
|U_{t+1}|=|U_t|-|S_{i_t}\cap U_t|\leq|U_t|-1\\
&\Longrightarrow\ \exists r\leq|U|:\ U_r=\varnothing.
\end{aligned}&&
\end{flalign*}
\begin{flalign*}
&\text{2. Feasibility:}\quad
\mathcal A(I)\text{ is feasible}\ \iff\
\mathcal A(I)\subseteq S\ \land\ \bigcup_{S_j\in\mathcal A(I)}S_j=U.&&
\end{flalign*}
\begin{flalign*}
&\begin{aligned}
&T_t:=\{S_{i_s}:0\leq s<t\},\qquad T_0=\varnothing,\qquad T_{t+1}=T_t\cup\{S_{i_t}\},\\
&T_t\subseteq S\ \land\ U_t=U\setminus\bigcup_{S_j\in T_t}S_j\qquad\text{(induction on }t\text{)},\\
&\mathcal A(I)=T_r\subseteq S\ \land\
\varnothing=U_r=U\setminus\bigcup_{S_j\in T_r}S_j
\ \Longrightarrow\ \bigcup_{S_j\in\mathcal A(I)}S_j=U.
\end{aligned}&&
\end{flalign*}
\begin{flalign*}
&\text{3. Efficiency:}\quad
\mathcal A\text{ is efficient}\ \iff\
\exists q\text{ polynomial}\ \forall I\in\mathcal I:\
\operatorname{time}_{\mathcal A}(I)\leq q(|I|).&&
\end{flalign*}
\begin{flalign*}
&\begin{aligned}
&\#\text{ operations per call}\in O(|S|\,|U|),\qquad
\#\text{ recursive calls}=r+1\leq|U|+1,\\
&|U|\leq|I|\ \land\ |S|\leq|I|\ \Longrightarrow\
\#\text{ operations}\in O(|S|\,|U|^2)\subseteq O(|I|^3),\\
&\operatorname{time}_{\text{one operation}}\in\poly(|I|)
\ \Longrightarrow\ \operatorname{time}_{\mathcal A}(I)\in\poly(|I|).
\end{aligned}&&
\end{flalign*}
\end{proof}
\end{minipage}\par

\noindent\begin{minipage}{\linewidth}
\begin{definition}[Harmonic number]
\[
H_0:=0,\qquad H_n:=\sum_{\ell=1}^{n}\frac1\ell
\qquad(n\in\mathbb N_{\geq1}).
\]
\[
H_n:\text{ the }n\text{-th harmonic number}.
\]
\end{definition}
\end{minipage}\par

\noindent\begin{minipage}{\linewidth}
\begin{lemma}[Price of the first greedy set]\label{lem:greedy-price}
\begin{flalign*}
&\begin{aligned}
&I=(U,S,w)\in\mathcal I,\qquad n:=|U|,\\
&\OPT^*:=\OPT(I),\qquad \OPT:=\OPT^*\setminus\{\varnothing\},\qquad c(\OPT)=c(\OPT^*),\\
&S_i\in\operatorname*{arg\,min}_{S_j\in S}\ppi(S_j)
\quad\Longrightarrow\quad \ppi(S_i)\leq\frac{c(\OPT)}{n}.
\end{aligned}&&
\end{flalign*}
\end{lemma}
\begin{proof}
\setlength{\jot}{2pt}
\setlength{\abovedisplayskip}{3pt}
\setlength{\belowdisplayskip}{3pt}
\setlength{\abovedisplayshortskip}{3pt}
\setlength{\belowdisplayshortskip}{3pt}
\begin{flalign*}
&\begin{aligned}
&\bigcup_{O_j\in \OPT}O_j=U\ \Longrightarrow\ \sum_{O_j\in \OPT}|O_j|\geq n.\\
&c(\OPT)=\sum_{O_j\in \OPT}w(O_j)=\sum_{O_j\in \OPT}|O_j|\ppi(O_j).
\end{aligned}&&
\end{flalign*}
\begingroup\small
\setlength{\abovedisplayskip}{3pt}
\setlength{\belowdisplayskip}{3pt}
\setlength{\abovedisplayshortskip}{3pt}
\setlength{\belowdisplayshortskip}{3pt}
\begin{flalign*}
&\begin{aligned}
&\ppi(S_i)
=\min_{S_j\in S}\ppi(S_j)
\leq\min_{O_j\in \OPT}\ppi(O_j)
\leq\frac{\sum_{O_j\in \OPT}|O_j|\ppi(O_j)}{\sum_{O_j\in \OPT}|O_j|}
=\frac{c(\OPT)}{\sum_{O_j\in \OPT}|O_j|}
\leq\frac{c(\OPT)}{n}.
\end{aligned}&&
\end{flalign*}
\endgroup
\end{proof}
\end{minipage}\par

\noindent\begin{minipage}{\linewidth}
\begin{theorem}[Approximation ratio of greedy set cover]
\[
\mathsf{GreedySetCover}\text{ is an }H_n\text{-approximation algorithm}.
\]
\[
n:=|U|.
\]
\end{theorem}
\begin{proof}
\setlength{\jot}{2pt}
\setlength{\abovedisplayskip}{3pt}
\setlength{\belowdisplayskip}{3pt}
\setlength{\abovedisplayshortskip}{3pt}
\setlength{\belowdisplayshortskip}{3pt}
\begingroup\raggedright
\noindent\(\mathsf{GreedySetCover}\) is an \(H_n\)-approximation algorithm
\(\overset{\text{Def.~\ref{def:alpha-approximation}, Thm.~\ref{thm:greedy-efficient}}}{\iff}
\forall I\in\mathcal I:\)
\mbox{\(c\bigl(\mathcal A(I)\bigr)
\leq H_n\,c\bigl(\OPT(I)\bigr)\)}.\par
\endgroup
\begin{flalign*}
&\begin{aligned}
&\text{Fix arbitrary }I=(U,S,w)\in\mathcal I.\\[3pt]
&\text{Strong induction on }n.\\[3pt]
&\text{1. Base case: }n=1\ \Longrightarrow\ \forall S_j\in S:\ S_j\neq\varnothing\ \Longrightarrow\ S_j=U,\\
&c\bigl(\mathcal A(I)\bigr)=\min_{S_j\in S:\ S_j\neq\varnothing}w(S_j)
=c\bigl(\OPT(I)\bigr)
=H_1\,c\bigl(\OPT(I)\bigr).
\end{aligned}&&
\end{flalign*}
\begin{flalign*}
&\begin{aligned}
&\text{2. Inductive step: }n\geq2.\\[3pt]
&S_i\in\operatorname*{arg\,min}_{S_j\in S}\ppi(S_j)\qquad\text{(first greedy set)},\\
\end{aligned}&&
\end{flalign*}

\begin{flalign*}
&\begin{aligned}
&k:=|S_i|\geq1,\qquad U':=U\setminus S_i,\qquad |U'|=n-k,\\
&S':=\{\,S_j':=S_j\setminus S_i:S_j\in S\,\},\qquad
w(S_j'):=w(S_j),\qquad I':=(U',S',w),\\
&C':=c\bigl(\mathsf{GreedySetCover}(U',S',w)\bigr).
\end{aligned}&&
\end{flalign*}
\begin{flalign*}
&\begin{aligned}
&c\bigl(\mathcal A(I)\bigr)
=w(S_i)+C'\\
&\overset{\text{Def. of }\ppi}{=}k\,\ppi(S_i)+C'\\
&\leq k\,\ppi(S_i)+H_{n-k}c(\OPT)\\
&\overset{\text{Lemma~\ref{lem:greedy-price}}}{\leq} k\,\frac{c(\OPT)}{n}+H_{n-k}c(\OPT)
=\left(\frac{k}{n}+H_{n-k}\right)c(\OPT)
\leq H_nc(\OPT).
\end{aligned}&&
\end{flalign*}
\begin{flalign*}
&\begin{aligned}
&n=k\ \Longrightarrow\ C'=0,\\
&n>k:\quad \underbrace{C'\leq H_{n-k}c(\OPT(I'))}_{\mathrm{IH}}
\leq H_{n-k}c(\OPT).
\end{aligned}&&
\end{flalign*}
\begin{flalign*}
&\frac{k}{n}=\sum_{\ell=n-k+1}^{n}\frac1n
\leq\sum_{\ell=n-k+1}^{n}\frac1\ell
=H_n-H_{n-k}.&&
\end{flalign*}

\begin{flalign*}
&\mathrm{IH}:\text{ induction hypothesis}.&&
\end{flalign*}
\end{proof}
\end{minipage}\par


\noindent\begin{minipage}{\linewidth}
\begin{example}[Tight bound for greedy set cover]
\setlength{\jot}{2pt}
\setlength{\abovedisplayskip}{3pt}
\setlength{\belowdisplayskip}{3pt}
\setlength{\abovedisplayshortskip}{3pt}
\setlength{\belowdisplayshortskip}{3pt}
\begin{flalign*}
&\begin{aligned}
&k\in\mathbb N_{\geq2},\qquad
\forall i\in\{1,\ldots,k\}:\quad |A_i|=|B_i|=2^{i-1},\qquad S_i:=A_i\cup B_i,\\
&\forall i:\ A_i\cap B_i=\varnothing,\qquad
\forall i\neq j:\ S_i\cap S_j=\varnothing,\\
&U:=\bigcup_{i=1}^{k}S_i,\qquad
n:=|U|=\sum_{i=1}^{k}2^i=2(2^k-1).
\end{aligned}&&
\end{flalign*}
\begin{flalign*}
&\begin{aligned}
&S_{k+1}:=\bigcup_{i=1}^{k}A_i,\qquad S_{k+2}:=\bigcup_{i=1}^{k}B_i,\\
&S_{k+1}\cap S_{k+2}=\varnothing\ \land\ S_{k+1}\cup S_{k+2}=U,\\
&S:=\{S_1,\ldots,S_{k+2}\},\qquad
\forall S_i\in S:\ w(S_i):=1,\qquad I_k:=(U,S,w).
\end{aligned}&&
\end{flalign*}
\begin{flalign*}
&\begin{aligned}
&j\in\{1,\ldots,k\},\qquad U_j:=\bigcup_{i=1}^{j}S_i,\qquad |U_j|=2(2^j-1),\\
&|S_j\cap U_j|=2^j,\qquad
|S_{k+1}\cap U_j|=|S_{k+2}\cap U_j|=\sum_{i=1}^{j}2^{i-1}=2^j-1,\\
&\forall i\leq j:\quad
\ppi(S_j\cap U_j)=\frac1{2^j}\leq\frac1{2^i}=\ppi(S_i\cap U_j),\\
&\ppi(S_j\cap U_j)=\frac1{2^j}<\frac1{2^j-1}
=\ppi(S_{k+1}\cap U_j)=\ppi(S_{k+2}\cap U_j),\\
&\forall i>j,\ i\leq k:\quad S_i\cap U_j=\varnothing\ \Longrightarrow\ \ppi(S_i\cap U_j)=\infty.
\end{aligned}&&
\end{flalign*}
\begin{flalign*}
&\begin{aligned}
&S_k\longrightarrow S_{k-1}\longrightarrow\cdots\longrightarrow S_1,\\
&\mathcal A(I_k)=\{S_1,\ldots,S_k\},\qquad c\bigl(\mathcal A(I_k)\bigr)=k.
\end{aligned}&&
\end{flalign*}
\begin{flalign*}
&\begin{aligned}
&\max_{S_i\in S}|S_i|=2^k<2(2^k-1)=n
\ \Longrightarrow\ c\bigl(\OPT(I_k)\bigr)\geq2,\\
&S_{k+1}\cup S_{k+2}=U\ \land\ c(\{S_{k+1},S_{k+2}\})=2
\ \Longrightarrow\ c\bigl(\OPT(I_k)\bigr)=2.
\end{aligned}&&
\end{flalign*}
\begin{flalign*}
&\frac{c\bigl(\mathcal A(I_k)\bigr)}{c\bigl(\OPT(I_k)\bigr)}
=\frac{k}{2}=\frac12\log_2\!\left(\frac n2+1\right)
=\Theta(\log n)=\Theta(H_n).&&
\end{flalign*}
\end{example}
\end{minipage}\par


\noindent\begin{minipage}{\linewidth}
\begin{theorem}[Hardness of approximating set cover (Dinur--Steurer, 2014)]
\begin{flalign*}
&\forall\varepsilon\in(0,1):\quad
(1-\varepsilon)\ln|U|\text{-approximation of }\Pi_{\text{minimum set cover}}\text{ is NP-hard}.&&
\end{flalign*}
\end{theorem}
\begin{proof}
\end{proof}
\end{minipage}\par

\subsection*{Questions}
\addcontentsline{toc}{subsection}{Questions}
\begin{enumerate}
\item Should Definition~1.1 explicitly define $\alpha:\mathcal I\to\mathbb R_{\geq1}$
and use $\alpha(I)$ in the approximation inequalities, allowing instance-dependent
factors such as $\alpha(I)=H_{|U|}$ and constant factors as constant functions
$\alpha(I)=a$ for all $I$?
\item Should an $\alpha$-approximation algorithm be required to be efficient,
i.e., run in polynomial time? The introduction of the lecture notes describes
approximation algorithms as efficient, but Definition~1.1 only states the
approximation inequality.
\item Should the minimum set cover problem require $|U|\geq1$ so that the
$H_n$-approximation guarantee always has $H_n\geq1$? The proof of
Theorem~1.3 in the lecture notes uses $n=0$ as its base case, but $H_0=0$,
whereas approximation factors are normally required to satisfy $\alpha\geq1$.
If the empty universe is allowed, how should the approximation factor be
stated? Does Definition~1.1 allow $\alpha$ to depend on $n$, with
$\alpha(0)=1$ and $\alpha(n)=H_n$ for $n\geq1$, or should the empty case be
handled separately from the approximation guarantee?
\end{enumerate}

\clearpage
\section{Approximation schemes}

\noindent\begin{minipage}{\linewidth}
\begin{definition}[Polynomial time approximation scheme (PTAS)]
\begin{flalign*}
&\begin{aligned}
&\mathcal A\text{ is a PTAS for }\Pi_{\mathrm{optimization}}\ \iff\\[3pt]
&\forall\varepsilon>0\ \exists q_\varepsilon\text{ polynomial}\ \forall I\in\mathcal I:\\[3pt]
&\quad\mathcal A_\varepsilon(I)\text{ terminates and is feasible for }I\ \land\\
&\quad\begin{cases}
c\bigl(\mathcal A_\varepsilon(I)\bigr)\leq(1+\varepsilon)c\bigl(\OPT(I)\bigr), & \text{minimization},\\[3pt]
c\bigl(\mathcal A_\varepsilon(I)\bigr)\geq(1-\varepsilon)c\bigl(\OPT(I)\bigr), & \text{maximization}
\end{cases}\ \land\\[3pt]
&\quad\mathrm{time}_{\mathcal A_\varepsilon}(I)\leq q_\varepsilon(|I|).
\end{aligned}&&
\end{flalign*}
\[
\Pi_{\mathrm{optimization}}:=(\mathcal I,\mathcal S),\qquad
\OPT(I)\in\mathcal S_I.
\]
\[
\forall\varepsilon>0:\quad\mathcal A_\varepsilon:\mathcal I\to\mathcal S.
\]

\[
\begin{aligned}
\varepsilon&:\text{ approximation parameter (fixed)},\\
\mathcal A_\varepsilon&:\text{ algorithm for fixed }\varepsilon,\\
q_\varepsilon&:\text{ polynomial in }|I|\text{ whose coefficients and degree may depend on }\varepsilon.
\end{aligned}
\]
\end{definition}
\end{minipage}\par

\noindent\begin{minipage}{\linewidth}
\begin{definition}[Fully polynomial time approximation scheme (FPTAS)]
\begin{flalign*}
&\begin{aligned}
&\mathcal A\text{ is an FPTAS for }\Pi_{\mathrm{optimization}}\ \iff\\[3pt]
&\exists q\text{ polynomial}\ \forall\varepsilon>0\ \forall I\in\mathcal I:\\[3pt]
&\quad\mathcal A(I,\varepsilon)\text{ terminates and is feasible for }I\ \land\\
&\quad\begin{cases}
c\bigl(\mathcal A(I,\varepsilon)\bigr)\leq(1+\varepsilon)c\bigl(\OPT(I)\bigr), & \text{minimization},\\[3pt]
c\bigl(\mathcal A(I,\varepsilon)\bigr)\geq(1-\varepsilon)c\bigl(\OPT(I)\bigr), & \text{maximization}
\end{cases}\ \land\\[3pt]
&\quad\mathrm{time}_{\mathcal A}(I,\varepsilon)\leq q\!\left(|I|,\frac1\varepsilon\right).
\end{aligned}&&
\end{flalign*}
\[
\Pi_{\mathrm{optimization}}:=(\mathcal I,\mathcal S),\qquad\OPT(I)\in\mathcal S_I.
\]
\[
\mathcal A:\mathcal I\times\mathbb R_{>0}\to\mathcal S.
\]
\[
\begin{aligned}
\varepsilon&:\text{ approximation parameter (input)},\\
q&:\text{ one fixed polynomial in }|I|\text{ and }1/\varepsilon.
\end{aligned}
\]
\end{definition}
\end{minipage}\par

\noindent\begin{minipage}{\linewidth}
\begin{example}[PTAS versus FPTAS running time]
\begin{flalign*}
&\begin{aligned}
&\mathrm{time}_{\mathcal A_\varepsilon}(I)\in O\!\left(|I|^{1/\varepsilon}\right).\\[3pt]
&\varepsilon>0\text{ fixed}:\quad |I|^{1/\varepsilon}\in\poly(|I|)
\quad\Longrightarrow\quad\text{PTAS running-time bound}.\\[3pt]
&|I|=2:\quad |I|^{1/\varepsilon}=2^{1/\varepsilon}\notin\poly(1/\varepsilon)
\quad\Longrightarrow\quad\text{not an FPTAS running-time bound}.
\end{aligned}&&
\end{flalign*}
\end{example}
\end{minipage}\par

\subsection*{Questions}
\addcontentsline{toc}{subsection}{Questions}
\begin{enumerate}
\item The introduction to Chapter~2 states that the greedy algorithms from
Chapter~1 approximate the optimum within a constant factor, including minimum
set cover. Is this a typo? The guarantee for greedy set cover is $H_n=\Theta(\log n)$,
where $n=|U|$, so its approximation factor is not constant. Should the statement
distinguish this logarithmic factor from the constant factors for matching
and vertex cover?
\item For maximization, should the PTAS and FPTAS guarantee be defined as
$c(\mathcal A(I,\varepsilon))\geq(1-\varepsilon)c(\OPT(I))$?
How does this relate to Definition~1.1, which uses
$c(\OPT(I))\leq\alpha\,c(\mathcal A(I))$?
Is the intended correspondence $\alpha=1/(1-\varepsilon)$ for
$0<\varepsilon<1$, rather than $\alpha=1+\varepsilon$?
\end{enumerate}

\clearpage
\section*{General questions}
\addcontentsline{toc}{section}{General questions}
\begin{enumerate}
\item When the lecture notes state a theorem but refer to a paper or another
source for its proof, are we expected to study and be able to reproduce that
proof, or is it sufficient to know the theorem's statement and how to use it?
\end{enumerate}

\end{document}
